\section[Nonradial Cubics and Pullbacks]{Nonradial Cubic Examples and Cubic Pullbacks}
\label{sec:nonradial-examples}

Cartan deformations attain the five-generator bound of
\cref{thm:harmonic-cubic-five-bound} in the dimensions allowed by
\cref{cor:nr-cartan-dimensions}. Their two weight eigenvalues in
\eqref{eq:nonradial-minimal-equation} distinguish them from the radial
irreducible cubics known to Fox and Tkachev
\cite[Introduction and Section~2.1]{FoxTkachev24}.
The models in \(\R^8,\R^{14}\) are unstable; five-generator inequalities
prove strict area minimality in \(\R^{26}\). We then pull back the
minimizing cubic targets.

\subsection{Nonradial harmonic models}
\label{sec:nonradial-models}

\subsubsection{One-parameter deformations of Cartan cubics}

Use the notation \(J_d=\operatorname{Herm}_3(\mathbb K_d)\) from
\cref{sec:exceptional-four-generator}. For the three Cayley--Dickson
inclusions \(\mathbb K_d\subset\mathbb K_{2d}\), \(d=1,2,4\), the
involution \(a+b\ell\mapsto a-b\ell\) preserves multiplication
and conjugation by the Cayley--Dickson formula. Its entrywise extension
therefore preserves the Jordan product and defines an automorphism
\(\sigma\) of \(J_{2d}\) with fixed algebra \(J_d\).
Relative to the common diagonal Jordan frame, it refines each
non-diagonal Peirce space as
\[
 J_{ij}\cong\mathbb K_{2d}
       =\mathbb K_d\oplus\mathbb K_d\ell,
 \qquad J_{ij}^{+}\cong\mathbb K_d,\quad
 J_{ij}^{-}\cong\mathbb K_d\ell.
\]
Thus the fixed algebra remains simple of rank three, and its
non-diagonal Peirce spaces have half the dimensions of those of
\(J_{2d}\).
On \(V=(J_{2d})_0\), with inner product
\(\ip{x}{y}=\frac12\tr(x\bullet y)\), this gives the orthogonal splitting
\[
 V=V_+\mathbin{\perp}V_-,\qquad
 V_+=(J_d)_0,\qquad V_-=U\cong\mathbb K_d^3.
\]
The Cartan cubic \(N=\det_{J_{2d}}|_V\) is invariant under \(\sigma\).
Hence \(N(b+u)=N(b-u)\) for \(b\in V_+\), \(u\in V_-\).
Since $N$ is cubic and even in $u$, it has the unique bidegree decomposition
\[
 N(b+u)=A(b)+\mathsf B(b,u),\qquad
 \deg A=(3,0),\quad \deg\mathsf B=(1,2).
\]
We deform the mixed component by setting
\begin{equation}\label{eq:nonradial-pencil}
 F_t=A+t\mathsf B,\qquad t\in\R.
\end{equation}
The choice of these pairs makes \(A=\det_{J_d}|_{(J_d)_0}\) itself a
Cartan cubic; see \cite[Theorem~3.1.1 and
Sections~3.2.6, 6.4.4--6.4.5]{NadirashviliTkachevVladut14} and
\cite{TkachevJordan14}. Both \(N\) and \(A\) are therefore harmonic.
Since \(\Delta=\Delta_b+\Delta_u\) and \(\mathsf B\) is linear in \(b\),
\[
 0=\Delta N=\Delta_bA+\Delta_u\mathsf B,
 \qquad \Delta F_t=(1-t)\Delta_bA=0.
\]
Thus the entire family is harmonic.
With \(\mathbb K=\mathbb K_d\), this identifies \(V\) with
\[
 V_d=\operatorname{Herm}_3(\mathbb K)_0\oplus\mathbb K^3,\qquad
 \ip{(Z,u)}{(Z',u')}=\tfrac12\operatorname{Re}\tr(ZZ')
                         +\operatorname{Re}(u^*u').
\]
Thus \(\dim V_d=6d+2\). The half-trace metric here differs from the
trace metric on the complexification in
\cref{sec:exceptional-four-generator}. Set
\begin{equation}\label{eq:nonradial-invariants}
 \rho=\tfrac12\tr Z^2,\quad \nu=|u|^2,\quad
 A=\det Z,\quad \mathsf{B}=-u^*Zu,\quad W=4u^*Z^2u.
\end{equation}
Then \(R=\rho+\nu\). For the nonradial member \(F_{-1/3}\), the
dictionary with the intrinsic invariants
\eqref{eq:nr-five-invariants} is
\[
 (r,s,a,p,h)=\left(\rho,\nu,\frac{3\sqrt3}{2}A,
             -\frac{\mathsf B}{2\sqrt3},\frac W{48}\right).
\]
The intrinsic normalization is \(\sqrt3 F_{-1/3}/4\), with
\(3/16\) times the quadratic weight below. We verify the model
by its coordinate table before invoking the intrinsic theorem.

To verify the mixed term and the isometry, write
\[
 Z=\begin{pmatrix}r_1&c&\bar b\\ \bar c&r_2&a\\ b&\bar a&r_3\end{pmatrix},
 \qquad r_1+r_2+r_3=0.
\]
In the determinant formula
\(r_1r_2r_3-r_1|a|^2-r_2|b|^2-r_3|c|^2+2\operatorname{Re}(abc)\),
replace \((a,b,c)\) by
\((a+v_1\ell,b+v_2\ell,c+v_3\ell)\). The quadratic part in
\(v\), with \(u=(\bar v_1,\bar v_2,\bar v_3)\), is
\(-u^*Zu=\mathsf{B}\). This verifies the isometry and the signs, including
the inclusion \(\mathbb H\subset\mathbb O\). In particular,
\[
 F_{-1/3}=\frac{4A-N}{3}
   =\frac13\tr\bigl(\sigma(x)\bullet x^{\bullet2}\bigr),
\]
where \(x\in(J_{2d})_0\) corresponds to \((Z,u)\).
After the common scaling above, this is the Cartan reconstruction in
\cref{cor:nr-cartan-dimensions}. These models realize all three ambient
dimensions allowed there; a classification of compatible splittings
would still be needed to classify all nonradial cubics.

\begin{proposition}[The nonradial differential algebra]
\label{prop:nonradial-algebra}
The algebra \(\R[\rho,\nu,A,\mathsf B,W]\) is Laplacian. For
\(F=F_{-1/3}\), it is the smallest Laplacian algebra containing
\(R,F\):
\[
 \cA_R(F)=\R[\rho,\nu,A,\mathsf{B},W].
\]
Its five generators are algebraically independent, and its spherical
quotient dimension is four. Their Laplacians and gradient products are
\begin{equation}\label{eq:nonradial-laplacians}
 \begin{gathered}
 \Delta\rho=6d+4,\qquad\Delta\nu=6d,\qquad
 \Delta A=\Delta \mathsf{B}=0,\\
 \Delta W=16d\rho+\tfrac{16}{3}(3d+2)\nu.
 \end{gathered}
\end{equation}
and, using symmetry, the following relations:
\begin{align}
 \Gamma(\rho,\rho)&=4\rho,&\Gamma(\rho,\nu)&=0,
       &\Gamma(\nu,\nu)&=4\nu,\notag\\
 \Gamma(\rho,A)&=6A,&\Gamma(\nu,A)&=0,
       &\Gamma(\rho,\mathsf{B})&=2\mathsf{B},\notag\\
 \Gamma(\nu,\mathsf{B})&=4\mathsf{B},&\Gamma(\rho,W)&=4W,
       &\Gamma(\nu,W)&=4W.\label{eq:nonradial-quadratic-pairings}
\end{align}
\begin{equation}\label{eq:nonradial-closure}
\begin{gathered}
 \begin{aligned}
 \Gamma(A,A)&=\tfrac43\rho^2,
 &\qquad\Gamma(\mathsf{B},\mathsf{B})&=\tfrac43\nu^2+W,\\
 \Gamma(A,\mathsf{B})&=\tfrac43\rho\nu-\tfrac12W,
 &\qquad\Gamma(A,W)&=16\nu A-\tfrac{16}{3}\rho \mathsf{B},
 \end{aligned}\\
 \Gamma(\mathsf{B},W)=16\rho \mathsf{B}-16\nu A+\tfrac{32}{3}\nu \mathsf{B},\\
 \Gamma(W,W)=16(\rho+\nu)W+\tfrac{64}{3}\mathsf{B}^2-64A\mathsf{B}.
\end{gathered}
\end{equation}
\end{proposition}

\begin{proof}
Put \(Q=uu^*\) and \(m_j=u^*Z^ju\). The trace-free characteristic
identity and the relevant gradients are
\begin{gather*}
 Z^3=\rho Z+AI,\qquad
 \nabla_ZA=2(Z^2-\tfrac23\rho I),\\
 \nabla_Z\mathsf{B}=-2(Q-\tfrac13\nu I),\qquad \nabla_u\mathsf{B}=-2Zu,\\
 \nabla_ZW=8(ZQ+QZ-\tfrac23m_1I),\qquad \nabla_uW=8Z^2u.
\end{gather*}
Their pairings give \eqref{eq:nonradial-closure}, including
\[
 |\nabla_ZW|^2=16\nu W+\tfrac{64}{3}\mathsf{B}^2,\qquad
 |\nabla_uW|^2=64m_4=16\rho W-64A\mathsf{B}.
\]
An orthonormal basis \(E_\alpha\) of \(\operatorname{Herm}_3(\mathbb K)_0\),
with the indicated trace metric, satisfies
\(\sum_\alpha E_\alpha^2=\frac23(3d+2)I\).
Taking diagonal and off-diagonal basis elements proves this identity;
it gives the last formula in \eqref{eq:nonradial-laplacians}.
The other Laplacians follow directly by taking traces. The chain
rule therefore proves that the displayed algebra is Laplacian.

For \(F=F_{-1/3}\) and \(S=|\nabla F|^2\), the table gives
\begin{equation}\label{eq:nonradial-gradient-square}
 S=\tfrac43(\rho-\nu/3)^2+\tfrac49W,\qquad
 \Delta S=\tfrac{32}{27}\bigl((15d+18)\rho+(3d+2)\nu\bigr).
\end{equation}
It also gives
\[
 \Gamma(F,S)=16(\rho+7\nu/27)F.
\]
Thus \(\cL(F)=-qF\), where \(q=8(\rho+7\nu/27)\) is positive
and nonradial. We check the remaining hypotheses needed for
\cref{thm:nonradial-five-generators}.

The Cartan cubic \(A\) is irreducible by its harmonic eikonal identity
\(|\nabla A|^2=4\rho^2/3\). Indeed, if a harmonic cubic
\(P\) in \(n>2\) variables satisfying
\(|\nabla P|^2=\kappa|x|^4\), \(\kappa>0\), had a real linear
factor, write it as
\(P(t,y)=t(at^2+t\ell(y)+Q(y))\).
Restriction to \(t=0\) gives \(Q=b|y|^2\), \(b^2=\kappa\).
Harmonicity gives \(\ell=0\) and \(6a+2b(n-1)=0\).
The coefficient of \(t^2|y|^2\) in the eikonal equation then
gives \(6ab+4b^2=2\kappa\), forcing \(n=2\), a contradiction.
Any real factorization of \(F\) would have a linear factor;
restriction to \(u=0\) would factor \(A\). Thus \(F\) is
irreducible.

If \(D_vF\equiv0\) for a constant vector \(v\), differentiating
\(\cL(F)=-qF\) gives \((D_vq)F=0\). Since \(q\) is positive
definite, this forces \(v=0\). The point \((Z,u)=(\operatorname{diag}(1,-1,0),0)\) is a
regular zero, and the two distinct eigenvalues of \(q\) establish
nonradiality.

The coordinate dictionary above now transfers
\cref{thm:nonradial-five-generators}, with \(m=3d\), to these
models, proving the smallest-algebra identity, algebraic independence,
and spherical quotient dimension four.
\end{proof}

\begin{theorem}[Minimal members of the Cartan deformation]
\label{thm:nonradial-models}
For \(d=1,2,4\), every cubic \(F_t\) in
\eqref{eq:nonradial-pencil} is harmonic. Its regular zero set is
minimal if and only if
\[
 t\in\{0,1,-\tfrac13\}.
\]
The members \(F_0=A\) and \(F_1=N\) are, respectively, a cylindrical
extension of a Cartan cubic and the original Cartan cubic.
The cubic \(F=F_{-1/3}\) is irreducible over \(\R\), has no
cylindrical directions, and satisfies
\begin{equation}\label{eq:nonradial-minimal-equation}
 \Delta F=0,\qquad
 \Gamma(F,S)=2qF,\qquad
 \cL(F)=-qF,\qquad q=8(\rho+7\nu/27).
\end{equation}
Its regular zero set is a minimal hypersurface in ambient dimension
\(8,14,26\), respectively. Its quadratic weight has two distinct
positive eigenvalues, with ratio \(7/27\).
\end{theorem}

\begin{proof}
Harmonicity was proved above. The closure table gives
\begin{equation}\label{eq:nonradial-pencil-gradient}
 S_t:=|\nabla F_t|^2
   =\tfrac43(\rho+t\nu)^2+t(t-1)W.
\end{equation}
Pairing its gradient with \(\nabla F_t\) yields
\begin{equation}\label{eq:nonradial-pencil-minimality}
 \cL(F_t)=-8\bigl(\rho+t^2(2-t)\nu\bigr)F_t
 -\tfrac83t(t-1)(3t+1)(\rho+t\nu)\mathsf B.
\end{equation}
This proves minimality for the three stated values and gives
\eqref{eq:nonradial-minimal-equation} at \(t=-1/3\).
For necessity, take \(Z=\operatorname{diag}(1,1,-2)\) and, for
\(t\ne0\), choose
\[
 u=t^{-1/2}e_3\quad(t>0),\qquad
 u=\sqrt{-2/t}\,e_1\quad(t<0).
\]
Here \(e_i\) are the standard coordinate vectors of \(\mathbb K^3\).
Then \(A=-2\), \(\mathsf B=2/t\), and \(F_t=0\), whereas
\(\nabla_uF_t=-2tZu\ne0\). Moreover, \(\rho+t\nu\) equals
\(4\) or \(1\), respectively. Thus
\eqref{eq:nonradial-pencil-minimality} cannot vanish at this regular
zero unless \(t=1\) or \(t=-1/3\).

The remaining assertions for \(F=F_{-1/3}\) were established in
\cref{prop:nonradial-algebra} while verifying the hypotheses of the
intrinsic five-generator theorem.
\end{proof}

These cones are not congruent to radial cubic cones. Indeed, if
\(P\circ Q\) has the same zero set as \(F\), with \(Q\) orthogonal
and \(P\) a radial cubic, the smooth-patch divisibility fact proved
in \cref{prop:minimizing-irreducibility} gives \(P\circ Q=cF\).
Here \(F\) is irreducible and has regular real zeros by
\cref{lem:square-free-critical-set}. Orthogonal equivariance of
\(\cL\) and \(\cL(cF)=c^3\cL(F)\) would then make \(c^2q\)
a multiple of \(R\), contradicting
\eqref{eq:nonradial-minimal-equation}.

For the rest of this section, put \(F=F_{-1/3}\) and
\(S=|\nabla F|^2\). The generator degrees give
\[
 (\cA_R(F))_2=\operatorname{span}\{\rho,\nu\},\qquad
 (\cA_R(F))_3=\operatorname{span}\{A,\mathsf{B}\}.
\]
Thus the cubic part is two-dimensional, unlike the hypothesis of
\cref{thm:cubic-one-dimensional-degree-three}.

\subsubsection{Singularities and the two unstable models}
\label{sec:nonradial-singularities}

\begin{lemma}\label{lem:nonradial-singularities}
The singular set of \(\{F=0\}\), equivalently the critical set of \(F\), is
\begin{equation}\label{eq:nonradial-critical-set}
 \Sing\{F=0\}=\Sigma=\{(Z,u):\nu=3\rho,\ Zu=0\}.
\end{equation}
The punctured set \(\Sigma\setminus\{0\}\) is smooth of
dimension \(4d\). At
\[
 c=(\operatorname{diag}(1,-1,0),(0,0,\sqrt3)),
\]
the Hessian of \(F\) has eigenvalues \(8/3,-8/3,0\), with
multiplicities \(d+1,d+1,4d\), respectively, and its kernel is
\(T_c\Sigma\). Every nonzero point of \(\Sigma\) is obtained
from \(c\) by a Hermitian change of basis and a positive dilation.
The zero sets, rescaled about such a point, converge smoothly on
compact subsets of the regular part of
\begin{equation}\label{eq:nonradial-transverse-cone}
 \R^{4d}\times\{(v_+,v_-)\in\R^{d+1}\oplus\R^{d+1}:
                           |v_+|=|v_-|\}.
\end{equation}
\end{lemma}

\begin{proof}
Equation \eqref{eq:nonradial-gradient-square} proves the description
of \(\Sigma\). For a nonzero \(u\), the restriction of \(Z\)
to \(u^\perp\) is a trace-free Hermitian \(2\times2\) matrix
of a fixed positive norm. These matrices form a sphere of dimension
\(d\), so \(\Sigma\setminus\{0\}\) is a smooth sphere bundle
over \(\mathbb K^3\setminus\{0\}\), of dimension \(4d\).
Its eigenvalues are \((\lambda,-\lambda,0)\), \(\lambda>0\),
which also proves the assertion about changes of basis and dilations.

At \(c\), let the diagonal variation of \(Z\) be
\((\xi,\eta,-\xi-\eta)\), let \(h_i\) be the variations of
\(u_i\), and use \(a,b,c'\) for the off-diagonal variations
in the matrix displayed above. The quadratic term of \(F\) is
\[
 Q_2=\xi^2-\eta^2-\tfrac2{\sqrt3}(\xi+\eta)
             \operatorname{Re}h_3
       +|b+\bar h_1/\sqrt3|^2-|a-h_2/\sqrt3|^2.
\]
The first three terms become \(8st'/3\) under the orthonormal
changes
\[
 s=\tfrac{\sqrt3}{2}(\xi+\eta),\qquad
 t=\tfrac12(\xi-\eta),\qquad
 t'=\tfrac12(\sqrt3t-\operatorname{Re}h_3).
\]
Completing to orthonormal coordinates gives
\(Q_2=\frac43(|v_+|^2-|v_-|^2)\) and the stated spectrum.
Differentiating \(\nabla F=0\) along \(\Sigma\), and comparing
dimensions, identifies the kernel. Finally,
\(r^{-2}F(c+rx)=Q_2(x)+rF(x)\), so
\cref{sec:tangent-instability-transfer} gives the tangent and smooth
regular graphical convergence.
The transverse cone is not a hyperplane, so every nonzero point of
\(\Sigma\) is geometrically singular; the vertex is singular by homogeneity.
\end{proof}

By \cref{prop:harmonic-positive-multiplier-obstruction}, a smooth
nonnegative subcalibrating multiplier must vanish on
\(\Sigma\setminus\{0\}\); in particular, radial powers fail.

\begin{proposition}\label{prop:nonradial-unstable}
The nonradial cones in ambient dimensions \(8\) and \(14\)
are unstable on their regular parts and are not area minimizing.
\end{proposition}

\begin{proof}
The transverse cone in \eqref{eq:nonradial-transverse-cone} has
dimension \(2d+1\) and squared second fundamental form
\(2d/r^2\). Its radial second variation, up to a positive factor,
is
\[
 \int_0^\infty\left(|\phi'|^2-\frac{2d}{r^2}\phi^2\right)
                                                      r^{2d}\,dr.
\]
The sharp Hardy constant \((2d-1)^2/4\) is smaller than \(2d\)
for \(d=1,2\), so a compactly supported test function on a radial
annulus gives a negative value. \Cref{lem:tangent-instability-transfer},
using the regular graphical convergence in \cref{lem:nonradial-singularities},
proves instability of the original cone.
\end{proof}

The boundary formula suggests a gradient factor.
Put \(K=\Delta S\) and \(T=|\nabla S|^2\). For \(H=F\sqrt S\),
the boundary speed is \(w=S^{-1}\).
On \(C_{\reg}\), the closure table gives
\(T/S=\frac{64}{27}(9\rho+\nu)\). Using
\eqref{eq:nonradial-gradient-square} and
\eqref{eq:nonradial-minimal-equation} in
\cref{prop:harmonic-power-Jacobi} yields
\begin{equation}\label{eq:nonradial-Jacobi-trace}
 -\frac{J_C(S^{-1})}{S^{-1}}
 =\frac{K/2-2q-T/S}{S}
 =\frac{16(d-3)(5\rho+\nu)}{9S}.
\end{equation}
For \(d=1,2\), this candidate fails the boundary condition. For \(d=4\),
its trace is strictly positive; we now prove the global subcalibration sign.

\subsubsection{The minimizing nonradial cone in dimension twenty-six}
\label{sec:nonradial-certificate}

\begin{theorem}\label{thm:nonradial-minimizing}
For \(d=4\), both homogeneous potentials
\[
 H=F\sqrt S,\qquad G=F\sqrt{RS}
\]
generate two-sided subcalibrations for \(E=\{F<0\}\).
Their divergence signs are strict on \(\{F\ne0\}\).
Consequently \(E\) is locally perimeter minimizing and
\(\partial\llbracket E\rrbracket\) is strictly area minimizing,
with support exactly \(\{F=0\}\subset\R^{26}\).
\end{theorem}

\begin{proof}
We first prove the assertion for the degree-five potential \(H\).
The closure table gives
\begin{equation}\label{eq:nonradial-T}
 T=\tfrac{64}{27}(9\rho+\nu)S-\tfrac{1024}{81}F\mathsf{B}.
\end{equation}
On \(S>0\), differentiation yields
\begin{align}\label{eq:nonradial-H-derivatives}
 \Delta H&=\frac{F}{\sqrt S}
                  \left(2q+\tfrac12K-\frac{T}{4S}\right),\notag\\
 |\nabla H|^2&=S^2+2qF^2+\frac{F^2T}{4S}\ge S^2.
\end{align}
The resulting defect is
\begin{align}\label{eq:nonradial-intermediate}
 \frac{\cL(H)}{S^{3/2}F}
 ={}&\tfrac12K-2q-\frac TS-\frac{F\Gamma(F,q)}S\notag\\
 &+\frac{F^2}{S^2}
       \left(qK+2q^2-\tfrac12\Gamma(S,q)\right)
       -\frac{F\Gamma(F,T)}{8S^2}\notag\\
 &+\frac{F^2}{8S^3}
                    \left(TK-\tfrac12\Gamma(S,T)\right).
\end{align}
The quotient is initially written for \(F\ne0\); multiplying
through gives an identity on all of \(\{S>0\}\).
Substituting the table gives
\begin{equation}\label{eq:nonradial-certificate}
 \cL(H)=\frac{16F}{729S^{3/2}}\mathcal N,
\end{equation}
where
\begin{align}\label{eq:nonradial-N}
 \mathcal N={}&(405\rho+81\nu)S^3+144F\mathsf{B}S^2\notag\\
 &+F^2S(46692\rho^2+20040\rho\nu+2212\nu^2-4239S)\notag\\
 &-(4992\rho+896\nu)F^3\mathsf{B}-256\nu F^4.
\end{align}

To prove its sign, define
\[
\begin{aligned}
 S_* &=\tfrac43\rho^2+\tfrac{40}{27}\rho\nu+\tfrac4{27}\nu^2,\\
 K_0&=4992\rho+896\nu,\\
 H_0&=\tfrac89(36639\rho^2+12634\rho\nu+1555\nu^2).
\end{aligned}
\]
The invariant image satisfies
\begin{equation}\label{eq:nonradial-bounds}
 S\le S_*,\qquad 9F^2\le RS,\qquad
 16\mathsf{B}^2\le9\nu S,\qquad
 F\mathsf{B}\le\tfrac3{16}(9\rho+\nu)S.
\end{equation}
Indeed, the trace-free eigenvalues of \(Z\) obey
\(\lambda_i^2\le4\rho/3\), whence
\(W\le16\rho\nu/3\) and the first inequality follows from
\eqref{eq:nonradial-gradient-square}. The second is Euler's
identity and Cauchy--Schwarz. The third follows from
\(\mathsf{B}^2\le\nu u^*Z^2u=\nu W/4\le9\nu S/16\).
The last is \(T\ge0\) in \eqref{eq:nonradial-T}.
There is the exact decomposition
\begin{align}\label{eq:nonradial-positive-certificate}
 \mathcal N={}&405\rho S^3+36S^2(F+2\mathsf{B})^2
                      +S^2(81\nu S-144\mathsf{B}^2)+F^2SH_0\notag\\
 &+K_0F^2\left(\tfrac3{16}(9\rho+\nu)S-F\mathsf{B}\right)\notag\\
 &+256\nu F^2\left(\tfrac19RS-F^2\right)
                      +4275F^2S(S_*-S).
\end{align}
Expansion recovers \eqref{eq:nonradial-N}; every summand is
nonnegative by \eqref{eq:nonradial-bounds}. Since $F\ne0$ implies
$\rho>0$ and $S>0$, the first summand proves the required strict
sign for \(\cL(H)\) off \(\{F=0\}\).

For the degree-six potential \(G=\sqrt R\,H\), we use
\eqref{eq:nonradial-H-derivatives} and
\eqref{eq:nonradial-T}, which give
\[
 \Delta H=\frac F{\sqrt S}
 \left(\frac{512}{9}\rho+\frac{320}{27}\nu
                         +\frac{256F\mathsf{B}}{81S}\right).
\]
The second and third bounds in \eqref{eq:nonradial-bounds} imply
\[
 \frac{|F\mathsf{B}|}{S}\le\frac14\sqrt{\nu(\rho+\nu)}
                   \le\frac{\rho+2\nu}{8}.
\]
Thus the parenthesis is at least
\((4576\rho+896\nu)/81>0\) for \(R>0\), and
\(H\Delta H\ge0\). On \(\{S>0\}\), the degree-five
potential \(H\) is smooth, satisfies \(H\cL(H)\ge0\), and
has nonvanishing gradient by \eqref{eq:nonradial-H-derivatives}.
\Cref{prop:cubic-radial-lift}, with degree \(5\) and ambient
dimension \(26\ge2\cdot5+2\), therefore gives
\(G\cL(G)\ge0\), strictly where \(F\ne0\), as well as
\(\nabla G\ne0\) on \(\{S>0\}\).

The exceptional set has codimension ten by
\cref{lem:nonradial-singularities}. Both potentials extend continuously
by zero across it with the signs of \(F\), and their normalized gradients
agree with the outward normal at regular zeros. Thus
\cref{thm:removable-subcalibration} applies; harmonicity gives the full
zero-set support. For both fields the strict sign on the two nonempty
phases gives \eqref{eq:two-phase-flux}, so
\cref{thm:strict-minimality-flux} proves strict area minimality.
\end{proof}

The degree-five potential \(H\) also proves strict stability directly.
Indeed, the cone has dimension \(n=25\), and its subcalibration sign
gives a nonnegative Jacobi boundary defect. Hence
\cref{thm:ground-state-gap} yields
\[
 Q_C(u)\geq\left(5-\frac{25}{2}\right)^2\int_Cr^{-2}u^2
       =\frac{225}{4}\int_Cr^{-2}u^2,
 \qquad u\in C_c^\infty(C_{\reg}\setminus\{0\}).
\]


\subsection{Cubic targets}

The common cubic potential gives the following pullback construction,
including for targets with nonisolated singularities.

\begin{corollary}[Pullback of the common cubic potential]
\label{cor:common-cubic-pullback}
Let \(f\not\equiv0\) be a harmonic cubic on \(\R^m\), and put
\(q=|y|^2\), \(S=|\nabla_yf|^2\), and \(H=f\sqrt{qS}\).
Suppose, on \(S>0\), that
\[
 f\Gamma_y(f,S)\geq0,\qquad H\cL_y(H)\geq0.
\]
For every quadratic Clifford map \(P:\R^N\to\R^m\) with
\(N\geq32\), set \(F=f(P)\), \(Q=q(P)\), and \(T=S(P)\).
Then
\[
 G=R^2F\sqrt{QT}
\]
is a subcalibration potential for \(\{F<0\}\). Its boundary current
is locally mass minimizing and strictly area minimizing, with support
\(\{F=0\}\).
The defining polynomial has degree six and the potential has
homogeneity degree sixteen.
\end{corollary}

\begin{proof}
Put \(H_0=f\sqrt S\). Bochner's identity and \(\Delta_y f=0\) give
\[
 H_0\Delta_yH_0
 =f^2\left(|\Hess f|^2-
       \frac{|\Hess f(\nabla f)|^2}{S}\right)
       +f\Gamma_y(f,S)\geq0.
\]
Since \(H_0\) has degree five, differentiating \(H=\sqrt q\,H_0\)
yields
\[
 H\Delta_yH=qH_0\Delta_yH_0+(m+9)H_0^2\geq0.
\]
Thus the ratio condition holds with \(M=0\). Euler's identity gives
\(\nabla H\ne0\) where \(H\ne0\); at a zero with \(S>0\),
\(\nabla H=\sqrt{qS}\nabla f\ne0\).
The exceptional set \(Z=\{\nabla f=0\}\) is conical and algebraic,
and Euler's identity gives \(Z\subset\{f=0\}\).
Its codimension is at least two by harmonicity and
\cref{lem:square-free-critical-set}.
Now apply \cref{thm:clifford-area-transfer} with \(\ell=6\),
\(a=2\), and \(M=0\): its bound is \(N\geq1+24+4+3=32\).
Harmonicity gives both signs near every target zero and hence the
support and strict-minimality assertions.
\end{proof}

The first hypothesis holds for the radial cubic models, where
\(\Gamma(f,S)\) is a positive multiple of \(qf\), and for the
nonradial models of \cref{thm:nonradial-models}, whose quadratic
minimality weight is positive. The second holds in the minimizing ranges of
\cref{thm:cubic-unified-sextic,thm:nonradial-minimizing}.
Since \(q,f,S\in\cA_q(f)\), this potential is a function of the
smallest target algebra. For an essential target with \(r\) minimal
homogeneous generators and a pullback with \(Q<R^2\) somewhere,
\cref{thm:smallest-clifford-pullback} gives exactly \(r+1\) source
generators. Thus the three-, four-, and five-generator cubic models
yield four-, five-, and six-generator pullback algebras whenever the
two displayed inequalities hold.

For the minimizing nonradial target in \cref{thm:nonradial-minimizing},
the degree-five potential \(H_0=f\sqrt S\) already satisfies
\(H_0\cL(H_0)\ge0\), and the Bochner identity above gives
\(H_0\Delta H_0\ge0\). Thus \cref{thm:clifford-area-transfer}, with
\(\ell=5\), \(a=1\), and \(M=0\), gives the degree-twelve potential
\(G=RF\sqrt T\) for every such Clifford map with \(N\ge28\).

For Cartan targets, a weighted inequality gives the lower-degree
pullback \(G=H(P)\).


\subsubsection{Cartan targets with a polynomial pullback potential}

Cartan's harmonic cubics on \(\R^m\), for \(m=14,26\), can be
normalized so that
\begin{equation}\label{eq:Cartan-target-normalization}
 \Delta f=0,\qquad |\nabla f|^2=9q^2,\qquad |f|^2\leq q^3,
 \qquad q=|y|^2.
\end{equation}
Their multiplicities are four and eight, respectively; see
\cite[Theorems~2.5 and~4.1]{Chi20}.  The algebra \(\R[q,f]\) is
Laplacian, with \(\Gamma(q,f)=6f\).

\begin{theorem}[Cartan cubic pullbacks]
\label{thm:cartan-pullback}
Let \(f\) be either normalized Cartan cubic in
\eqref{eq:Cartan-target-normalization}, and let
\(P:\R^N\to\R^m\) be a quadratic Clifford map.  Then
\[
 F=f\circ P,\qquad Q=|P|^2,\qquad G=QF
\]
give a perimeter-minimizing phase \(\{F<0\}\).
Its boundary current \(C\) is strictly area minimizing and is supported
on \(\{F=0\}\). It is strictly stable, with the explicit bound
\begin{equation}\label{eq:Cartan-pullback-Hardy}
 Q_C(u)\geq
 \left[\left(10-\frac{N-1}{2}\right)^2+8m-106\right]
 \int_Cr^{-2}u^2,
 \qquad u\in C_c^\infty(C_{\reg}\setminus\{0\}).
\end{equation}
The polynomial \(F\) has degree six and the
subcalibration potential \(G\) has degree ten.  The construction
exists for \((m,N)=(14,256)\) and \((26,8192)\), and for direct
sums of either Clifford system.
\end{theorem}

\begin{proof}
Use the degree-five target potential \(H=qf\) from
\cref{ex:harmonic-isoparametric-radial}. The target closure gives
\begin{align}
 |\nabla H|^2&=9q^4+16qf^2,\qquad
 \Delta H=2(m+6)f,\label{eq:Cartan-target-gradient}\\
 \cL(H)&=9(2m-24)q^4f+16(2m+1)qf^3.
 \label{eq:Cartan-target-defect}
\end{align}
Substituting the target closure in
\(\cL(H)=|\nabla H|^2\Delta H-\Gamma(H,|\nabla H|^2)/2\) gives
the second identity. Hence
\begin{equation}\label{eq:Cartan-target-weighted-defect}
 2q\cL(H)-5H|\nabla H|^2
 =q^2f\bigl(9(4m-53)q^3+16(4m-3)f^2\bigr).
\end{equation}
Both coefficients in parentheses are positive for \(m=14,26\).
After multiplication by \(H\), this is precisely
\eqref{eq:pullback-target-sign} for \(\kappa=1/2\).

The smallest source algebra is \(\cA_R(F)=\R[R,Q,F]\) by
\cref{thm:smallest-clifford-pullback}. Indeed, the target algebra is
\(\cA_q(f)=\R[q,f]\), and \(|\nabla f|^2=9q^2\) excludes a
nonzero cylindrical direction. For example,
\[
 \Delta Q=8mR,\quad \Delta F=0,\quad
 \Gamma(Q,F)=24RF,\quad |\nabla F|^2=36RQ^2.
\]
The target inequality therefore gives the required source sign by
\cref{prop:pullback-subcalibration}; the expanded gradient and defect
are recorded in \cref{app:Cartan-source-certificate}.

By \eqref{eq:Cartan-target-gradient}, \(\nabla H\ne0\) off the
target origin. The pullback results
\cref{prop:pullback-subcalibration,lem:Clifford-algebra-pullback}
therefore give comparison across \(P^{-1}(0)\), of codimension \(m\)
away from the origin. Surjectivity makes \(F\ne0\); harmonicity
identifies the full support and gives both phases by
\cref{rem:polynomial-finite-perimeter}. The strict sign in
\eqref{eq:Cartan-pullback-defect} then invokes
\cref{cor:polynomial-strict-minimality}.

For the improved Hardy bound, compute the boundary term. On the regular cone, \(a=Q\) and \(|\nabla F|^2=36RQ^2\), so
\eqref{eq:boundary-potential-algebraic} gives
\[
 \beta_G=\left.\frac{R\cK_F(Q)}{Q^3|\nabla F|^2}\right|_C
 =(8m-96)\frac{R^2}{Q}-10\geq8m-106>0.
\]
Apply \cref{thm:ground-state-gap} with \(\ell=10\) to obtain
\eqref{eq:Cartan-pullback-Hardy}.

Finally, the standard Clifford module dimensions give systems with
fourteen matrices on \(\R^{256}\) and twenty-six matrices on
\(\R^{8192}\); see \cite[Table~4.1]{RiedlerThesis25}.  Direct sums preserve
the Clifford relations and all the identities used above.
\end{proof}
